\subsection{Main Multi-Qubit Gates}\label{sec:main-multi-qubit-gates}

\subsubsection{Controlled NOT (CNOT) Gate}\label{sec:controlled-not-cnot-gate}

The \definition{Controlled NOT (CNOT) Gate} is a \textbf{two-qubit gate} composed of:
\begin{itemize}
    \item A \textbf{control qubit} that \hl{determines} whether the \hl{NOT operation is applied or not}.
    \item A \textbf{target qubit} that is \hl{flipped} (NOT operation) if the control qubit is in the state $\ket{1}$, and remains unchanged if the control qubit is in the state $\ket{0}$.
\end{itemize}
In simple terms, it is a gate that \hl{performs a NOT operation on the \emph{target} qubit only when the \emph{control} qubit is in the state $\ket{1}$.}

\highspace
\begin{flushleft}
    \textcolor{Green3}{\faIcon{book} \textbf{Action on Basis States}}
\end{flushleft}
The CNOT gate acts on the computational basis states as follows:

\begin{table}[!htp]
    \centering
    \begin{tabular}{@{} c | c @{}}
        \toprule
        \textbf{Input} & \textbf{Output} \\
        \midrule
        $\ket{00}$ & $\ket{00}$ \\
        $\ket{01}$ & $\ket{01}$ \\
        $\ket{10}$ & $\ket{11}$ \\
        $\ket{11}$ & $\ket{10}$ \\
        \bottomrule
    \end{tabular}
\end{table}

\noindent
If the first qubit (control) is $\ket{0}$, the second qubit (target) remains unchanged. If the first qubit is $\ket{1}$, the second qubit is flipped: $\ket{0} \leftrightarrow \ket{1}$.

\highspace
\begin{flushleft}
    \textcolor{Green3}{\faIcon{square-root-alt} \textbf{Matrix Representation}}
\end{flushleft}
The CNOT gate can be represented as a $4 \times 4$ unitary matrix in the computational basis $\{\ket{00}, \ket{01}, \ket{10}, \ket{11}\}$:
\begin{equation}
    \text{CNOT} = \begin{bmatrix}
        1 & 0 & 0 & 0 \\
        0 & 1 & 0 & 0 \\
        0 & 0 & 0 & 1 \\
        0 & 0 & 1 & 0
    \end{bmatrix}
\end{equation}
This matrix indicates that the first two basis states $\ket{00}$ and $\ket{01}$ are unchanged, while $\ket{10}$ and $\ket{11}$ are swapped, reflecting the action of the CNOT gate on the target qubit based on the state of the control qubit.

\highspace
\begin{flushleft}
    \textcolor{Green3}{\faIcon{book} \textbf{Action on General States}}
\end{flushleft}
For a general two-qubit state $\ket{v} = c_{0}\ket{00} + c_{1}\ket{01} + c_{2}\ket{10} + c_{3}\ket{11}$, the CNOT gate transforms it as follows:
\begin{equation}
    \text{CNOT} \ket{v} = c_{0}\ket{00} + c_{1}\ket{01} + c_{2}\ket{11} + c_{3}\ket{10}
\end{equation}
The amplitudes $c_{0}$ and $c_{1}$ remain unchanged, while $c_{2}$ and $c_{3}$ are \textbf{swapped}, demonstrating the conditional nature of the CNOT gate based on the state of the control qubit.

\highspace
\begin{flushleft}
    \textcolor{Green3}{\faIcon{tools} \textbf{Circuit Representation}}
\end{flushleft}
In quantum circuit diagrams, the CNOT gate is typically represented with a solid dot on the control qubit line and a plus sign $\oplus$ (indicating the NOT operation) on the target qubit line, connected by a vertical line.

\begin{center}
    \begin{quantikz}
        \lstick{$\ket{v_{A}}$} & \ctrl{1}   & \rstick{$\ket{v_{A}}$} \\
        \lstick{$\ket{v_{B}}$} & \targ{}    & \rstick{$\ket{v_{A}} \oplus \ket{v_{B}}$}
    \end{quantikz}
\end{center}

\noindent
Here, $\ket{v_{A}}$ is the control qubit and $\ket{v_{B}}$ is the target qubit. The output state of the target qubit is the result of the XOR operation between the control and target qubits, reflecting the conditional flipping of the target qubit based on the state of the control qubit.

\highspace
In summary, \hl{if the control qubit is $\ket{1}$, apply NOT (Pauli-X) to the target qubit; otherwise do nothing.}

\highspace
\begin{flushleft}
    \textcolor{Green3}{\faIcon{cogs} \textbf{Properties}}
\end{flushleft}
\begin{itemize}
    \item \important{Unitary}
    \begin{equation*}
        \text{CNOT}^{\dagger} \cdot \text{CNOT} = \begin{pmatrix}
            1 & 0 & 0 & 0 \\
            0 & 1 & 0 & 0 \\
            0 & 0 & 0 & 1 \\
            0 & 0 & 1 & 0
        \end{pmatrix}^{\dagger} \cdot \begin{pmatrix}
            1 & 0 & 0 & 0 \\
            0 & 1 & 0 & 0 \\
            0 & 0 & 0 & 1 \\
            0 & 0 & 1 & 0
        \end{pmatrix} = I
    \end{equation*}
    CNOT gate is unitary, so it is a valid quantum gate that preserves normalization of quantum states, and the total probability of all outcomes remains 1.

    \item \important{Reversible}
    \begin{equation*}
        \text{CNOT}^{-1} = \text{CNOT}^{\dagger} = \text{CNOT}
    \end{equation*}
    The CNOT gate is reversible, meaning that applying the CNOT gate twice will return the system to its original state:
    \begin{equation*}
        \text{CNOT} \left(\text{CNOT}\left(\ket{\psi}\right)\right) = \ket{\psi}
    \end{equation*}
    In other words, we can \textbf{always undo a CNOT operation without loss of information}.

    \item \important{Conditional Operation}
    \begin{equation*}
        \text{CNOT}\left(\ket{v_A} \otimes \ket{v_{B}}\right)
        \neq
        \text{CNOT}\ket{v_A} \otimes \text{CNOT}\ket{v_{B}}
    \end{equation*}
    The CNOT gate is a conditional operation, meaning that the \textbf{action on the target qubit depends on the state of the control qubit}. This is different from parallel gates, where each qubit is acted upon independently:
    \begin{equation*}
        \left(H \otimes X\right)\ket{00} = \left(H \ket{0} \right) \otimes \left( X \ket{0} \right)
    \end{equation*}
    CNOT \textbf{breaks parallelism} because it introduces interaction: the \hl{evolution of one qubit depends on another}, so the \textbf{system can no longer be described as independent parts}.
    
    \item \important{Entangling}. CNOT creates \emph{entanglement} \textbf{if and only if} the input state has superposition on the control qubit (in the computational basis).
    
    For example:
    \begin{equation*}
        \left(H \otimes I\right)\ket{00} = \frac{1}{\sqrt{2}} \left( \ket{00} + \ket{10} \right)
    \end{equation*}
    Applying CNOT to this state:
    \begin{equation*}
        \text{CNOT} \left( \frac{1}{\sqrt{2}} \left( \ket{00} + \ket{10} \right) \right) = \frac{1}{\sqrt{2}} \left( \ket{00} + \ket{11} \right)
    \end{equation*}
    The resulting state is an entangled state (cannot be factored into a tensor product, \autopageref{ex:not-all-states-can-be-factored}), demonstrating that CNOT can create entanglement from a superposition state of the control qubit.
\end{itemize}

\highspace
\begin{flushleft}
    \textcolor{Green3}{\faIcon{tools} \textbf{Multiple CNOT Gates}}
\end{flushleft}
Multi-CNOT notation means applying several CNOT gates that share the same control qubit, and these can be represented equivalently as sequential gates or as a single control connected to multiple targets. Visually, this can be represented as:
\begin{center}
    \begin{quantikz}
        \lstick{$\ket{v_{A}}$} & \ctrl{1}   & \ctrl{2}   & \rstick{$\ket{v_{A}}$} \\
        \lstick{$\ket{v_{B}}$} & \targ{}    & \qw        & \rstick{$\ket{v_{A}} \oplus \ket{v_{B}}$} \\
        \lstick{$\ket{v_{C}}$} & \qw        & \targ{}    & \rstick{$\ket{v_{A}} \oplus \ket{v_{C}}$}
    \end{quantikz}
    $\equiv$
    \begin{quantikz}
        \lstick{$\ket{v_{A}}$} & \ctrl{2}   & \ctrl{1}   & \rstick{$\ket{v_{A}}$} \\
        \lstick{$\ket{v_{B}}$} & \qw        & \targ{}    & \rstick{$\ket{v_{A}} \oplus \ket{v_{B}}$} \\
        \lstick{$\ket{v_{C}}$} & \targ{}    & \qw        & \rstick{$\ket{v_{A}} \oplus \ket{v_{C}}$}
    \end{quantikz}
\end{center}

\noindent
But it can also be represented as:
\begin{center}
    \begin{quantikz}
        \lstick{$\ket{v_{A}}$} & \ctrl{2}   & \rstick{$\ket{v_{A}}$} \\
        \lstick{$\ket{v_{B}}$} & \targ{}    & \rstick{$\ket{v_{A}} \oplus \ket{v_{B}}$} \\
        \lstick{$\ket{v_{C}}$} & \targ{}    & \rstick{$\ket{v_{A}} \oplus \ket{v_{C}}$}
    \end{quantikz}
\end{center}

\noindent
The \hl{order of the CNOT gates does not matter} because they \textbf{commute when they share the same control qubit}, meaning that the final state of the system will be the same regardless of the order in which the CNOT gates are applied.